% from ega4-2-fr.tex
\oldpage[IV-2]{28}
Comme le morphisme $Y'\to Y$ est fidèlement plat et quasi-compact (2.2.13 et
1.1.2), on peut, en vertu de (I, 3.3.11), se borner au cas où $Y=S$,
$Y'=S'$. Si $g'$ est la projection $X'\to X$, on sait que l'on a, pour toute
partie $M$ de $X$,
\[
g^{-1}(f(M))=f'(g'^{-1}(M))\quad\text{(I, 3.4.8).}
\]
Si $f'$ est un morphisme ouvert (resp. fermé) alors, pour toute partie
ouverte (resp. fermée) $M$ de $X$, $f'(g'^{-1}(M))$ est ouverte (resp.
fermée) dans $Y'$, et comme $g$ est fidèlement plat et quasi-compact, on en
conclut que $f(M)$ est ouverte (resp. fermée) dans $Y$ en vertu de (2.3.12).
Cela démontre la proposition dans les cas (i) et (ii). Démontrons-la dans le
cas (iii) (ce qui l'entraînera dans le cas (iv), compte tenu de (2.6.1,
(iv)). En vertu de (2.6.1, (ii)) $f$ est injectif, et il reste à prouver que
$f$, considéré comme application de $X$ sur $f(X)$, est une application
quasi-compacte et ouverte. Comme $f'$ est quasi-compact, il en est de même de
$f$ (1.1.4). Il suffit alors de prouver que pour toute partie fermée $Z$ de
$X$, on a
\[
Z=f^{-1}(\overline{f(Z)})~;
\]
comme $g'$ est surjective, cette relation équivaut à
\[
g'^{-1}(Z)=g'^{-1}(f^{-1}(\overline{f(Z)}))
=f'^{-1}(g^{-1}(\overline{f(Z)})).
\]
Or, comme $f$ est quasi-compact, il en est de même de sa composée avec
l'injection canonique $Z\to X$ ($Z$ étant ici le sous-préschéma fermé réduit
de $X$ ayant $Z$ comme espace sous-jacent). Appliquant (2.3.10) à la partie
$f(Z)$ de $Y$ (image du morphisme $f|Z$), on obtient
\[
g^{-1}(\overline{f(Z)})=\overline{g^{-1}(f(Z))}~;
\]
la formule à démontrer équivaut donc à la suivante~:
\[
Z'=f'^{-1}(\overline{f'(Z')}),
\]
où l'on a posé $Z'=g'^{-1}(Z)$~; mais cette formule résulte de l'hypothèse
que $f'$ est un homéomorphisme de $X'$ sur $f'(X')$.

\begin{remark}[2.6.3]
\label{IV.2.6.3-fr}
Dans les cas (i) et (ii), les conclusions de (2.6.2) sont encore valables
quand on suppose seulement $g$ quasi-fidèlement plat (2.3.3) et
quasi-compact~; en effet, en vertu de (2.1.4), (I, 3.5.2) et
(I, 9.1.13.1), on peut encore se ramener au cas où $Y=S$, $Y'=S'$~; la
conclusion résulte alors de (2.3.12). Dans les cas (iii) et (iv), les
conclusions restent valables quand on suppose seulement $g$ quasi-fidèlement
plat, pourvu que l'on suppose déjà $f$ quasi-compact~; en effet, on n'utilise
alors que (2.3.10) et le fait que $g$ est surjectif. Enfin, si $g$ est
fidèlement plat et localement de présentation finie ou si $g$ est surjectif
et $S$ discret, la conclusion de (2.6.2) est valable lorsque $P$ est la
propriété~:
\begin{enumerate}
\item[(iii \emph{bis})] être un homéomorphisme sur le sous-espace image~;
\end{enumerate}
cela résulte en effet de la démonstration donnée dans (2.6.2) et de la
Remarque (2.4.11).
\end{remark}

\begin{corollary}[2.6.4]
\label{IV.2.6.4-fr}
Les notations étant celles de (2.6.1), supposons le morphisme $g:S'\to S$
fidèlement plat et quasi-compact. Considérons pour un morphisme, la propriété
d'être~:
\begin{enumerate}
\item[(i)] universellement ouvert~;
\item[(ii)] universellement fermé~;
\item[(iii)] quasi-compact et universellement bicontinu~;
\item[(iv)] un homéomorphisme universel~;
\item[(v)] quasi-compact~;
\item[(vi)] quasi-compact et dominant.
\end{enumerate}
Alors, si $P$ désigne une des propriétés précédentes, pour que $f$ possède la
propriété $P$, il faut et il suffit que $f'$ la possède.
\end{corollary}

